\scanpage{18}
\deleted{\unsure{1|2}} A) Application à la comparaison avec les revêtements étales holomorphes

B) Application à \deleted{la construction} \inserted{une factorisation} canonique

\uline{Proposition} \quad Soit $f:X\longrightarrow Y$ un morphisme fini, \inserted{\uuline{univ. ouvert}}, \deleted{\unsure{sur}} avec $Y$ loc noethérien. \deleted{\unsure{Supposons|Pour}} \inserted{Pour} $y\in Y$, soit $d(y)$ le nb de pts géométriques de $f^{-1}(y)$ ($=$ degré séparable de $f^{-1}(y)$ sur $k(y)$]

\uline{Supposons} $d(y)$ \unsure{loc. cste.|constant} Alors $f$ se factorise en $X\xrightarrow{f_1}X'_0\xrightarrow{f_2}Y$ où $f_1$ est \inserted{fini,} surjectif radiciel, et $f_2$ est un morphisme \uline{fini} \uline{étale}. \inserted{\uline{Alors}} \uline{Tout} $Y$-morphisme $X\longrightarrow Z$, \margin{\unsure{1}} avec $Z$ étale sur $Y$, se factorise de façon unique en $X\longrightarrow X'\longrightarrow Z$.

On peut supposer \inserted{$Y$ connexe donc} $d(y)$ constant, soit $n$. (quitte à remplacer $Y$ par une partie à la fois ouverte et fermée convenable, pour \unsure{tt entier|la notion}]

% Mirrored printing shows through from the other face of the sheet.

% archive sheet number: 17
